\section{Relation to previous work}\label{sec:related}

The ingredients of Theorem~\ref{thm:central-obstruction} belong to several
established literatures.  Its content lies in the way their quantifiers meet.

\subsection{Heat-kernel lattice actions and small-time geometry}

Menotti and Onofri introduced the compact-group heat kernel as a plaquette
action and emphasized its weak-coupling and Hamiltonian features
\cite{MenottiOnofri1981}.  We use exactly this local action, with the heat
generator and metric fixed in Section~\ref{sec:setup}.  Our result is not a
new formula for the heat kernel.  It studies the product of incident
heat-kernel factors after conditioning on all links except one.

The logarithmic input is Varadhan's small-time identification of the heat
kernel with squared Riemannian distance
\cite{Varadhan1967Fundamental,Varadhan1967Diffusion}.  We need only the
uniform logarithmic form on a compact manifold, not a prefactor expansion.
This is why cut-locus multiplicity and the vanishing radial Haar factor in
$\SUtwo$ do not alter the theorem.  The new calculation after that input is
the weighted two-endpoint energy generated by an admissible incident-staple
configuration.

\subsection{Conditional criteria and product domination}

Dobrushin formulated uniqueness and comparison conditions in terms of the
sensitivity of single-site conditional distributions
\cite{Dobrushin1968}.  The coefficient convention in
\eqref{eq:dobrushin-coefficient} is the direct total-variation form of that
idea.  Theorem~\ref{thm:dobrushin-failure} does not challenge Dobrushin's
theorem: it gives explicit boundary pairs showing that its row hypothesis is
false for the microscopic link specification at sufficiently small heat
time.

Liggett, Schonmann, and Stacey proved product-domination results for dependent
binary fields under strong conditional-probability hypotheses
\cite{LiggettSchonmannStacey1997}.  If plaquette-good indicators were to be
controlled through a bound uniform over all external microscopic links, the
same quantifier would appear.  Equation~\eqref{eq:bad-limit} shows that this
particular route cannot begin from a vanishing worst-boundary bad probability.
It says nothing against domination theorems whose hypotheses are established
after blocking, on a restricted good boundary set, or under a boundary law
that retains the cost of the adverse configuration.

\subsection{Center twists, fluxes, and vortices}

Central boundary modifications are classical.  't~Hooft used twisted
boundary conditions to define electric and magnetic flux sectors
\cite{tHooft1979}; Mack and Petkova related changes of vorticity and boundary
conditions to a sufficient confinement mechanism
\cite{MackPetkova1980}; later lattice studies measured vortex or electric-flux
free energies, including the $\SUtwo$ analysis of de~Forcrand and von~Smekal
\cite{deForcrandVonSmekal2002}.

Those objects are global: they compare partition functions in topologically
or flux-distinguished sectors and, crucially, retain a free-energy cost.  The
boundary in Proposition~\ref{prop:central-surjectivity} is local, lives on an
open block, and carries no asserted topological sector.  We deliberately
condition on it and ask what a selected link does when the boundary cost is
treated as free.  Accordingly, the theorem neither discovers center
frustration nor computes a vortex free energy.  It demonstrates that erasing
such a cost by a worst-case conditional supremum reverses the desired local
tail estimate.

\subsection{Constructive renormalization}

Balaban's analysis of lattice gauge fields constructs localized effective
actions in a small-field approximation, develops cluster expansions, and
incorporates large-field regions into an ultraviolet-stability argument
\cite{Balaban1987,Balaban1988,Balaban1989}.  That machinery does not posit
that every microscopic boundary is a cost-free small field.  Our finite
conditional obstruction therefore does not displace, contradict, or close
the constructive renormalization problem.  Its relevance is diagnostic: any
proposed shortcut from a one-plaquette heat-kernel tail to an interacting
defect gas must specify where the action and entropy of a frustrating
boundary are paid.

\subsection{Contribution boundary}

The nearest established ingredients thus differ from the present theorem in
one of four ways: they define the heat-kernel action rather than its
frustrated conditional concentration; provide the heat-kernel asymptotic
rather than the lattice realization; formulate conditional criteria rather
than test them on this gauge specification; or retain global center-flux
free energies rather than take a local worst-boundary supremum.  We make no
unqualified priority claim.  The contribution established here is the
self-contained combination: an explicit admissible central staple pattern,
its compact-group metric barycenter law, an exact good-event exponent,
positive-measure stability, and the resulting link-coordinate Dobrushin row
obstruction.
